\section{Introduction}
\label{sec:introduction}

Sensitivity measures how many individual input bits can change a Boolean
function's value; block sensitivity allows disjoint sets of bits to be
flipped instead. We construct total Boolean functions for which the second
quantity grows faster than any constant multiple of the square of the first.

For a positive integer \(n\), write \([n]=\{1,\ldots,n\}\).
If \(x\in\{0,1\}^n\) and \(B\subseteq[n]\), let \(x^B\) be obtained
by flipping precisely the coordinates in \(B\); abbreviate
\(x^{\{a\}}\) to \(x^a\).
For a total Boolean function \(f:\{0,1\}^n\to\{0,1\}\), define
\[
 s(f,x)=\bigl|\{a\in[n]:f(x^a)\ne f(x)\}\bigr|,
 \qquad s(f)=\max_x s(f,x).
\]
Its block sensitivity \(\bs(f,x)\) is the largest number of pairwise
disjoint nonempty sets \(B_1,\ldots,B_m\subseteq[n]\) such that
\(f(x^{B_j})\ne f(x)\) for every \(j\), and
\(\bs(f)=\max_x\bs(f,x)\).

The Sensitivity Conjecture asked whether block sensitivity is bounded by a
polynomial in sensitivity, uniformly over the number of inputs and the
function. The question goes back to Nisan's work on parallel computation,
which introduced block sensitivity and already recorded Rubinstein's
quadratic separation \cite[Section~2]{Nisan1989}.
Nisan and Szegedy studied the problem through the degree of the unique
multilinear real polynomial representing a Boolean function, and explicitly
suggested the quadratic inequality \(\bs(f)\le s(f)^2\)
\cite[Section~4]{NisanSzegedy1994}.
Rubinstein's construction \cite{Rubinstein1995} showed that a polynomial
upper bound would need degree at least two. Virza improved the separation
to \(\bs(f)=\tfrac12s(f)^2+\tfrac12s(f)\) \cite{Virza2011}, and
Ambainis and Sun subsequently obtained examples satisfying
\[
 \bs(f)=\frac23s(f)^2-\frac13s(f)
\]
\cite[Theorem~1]{AmbainisSun2011}.
These improvements left open whether any universal quadratic upper bound
could hold.

Huang resolved the polynomial conjecture in 2019. His degree bound,
together with Tal's refinement of the degree--block-sensitivity comparison
\cite{Tal2013}, gives \(\bs(f)\le s(f)^4\) for every total Boolean function
\cite[Theorem~1.5]{Huang2019}. Wellens later improved its leading constant,
proving
\(\bs(f)\le\sqrt{2/3}\,s(f)^4+1\)
\cite[Corollary~4.2]{Wellens2022}.
These results settle the polynomial question. The theorem below rules out
its quadratic strengthening.

\begin{theorem}\label{thm:main}
For every real \(C>0\), there are a positive integer \(n\) and a
nonconstant total Boolean function \(f:\{0,1\}^n\to\{0,1\}\) such that
\[
 \bs(f)>C\,s(f)^2.
\]
More precisely, for every integer \(d\ge1\) the construction below gives
a nonconstant total Boolean function \(f\) with
\[
 \frac{\bs(f)}{s(f)^2}\ge \frac{2^d}{4(d+2)^2}.
\]
\end{theorem}

\paragraph{Method and relation to spectral sensitivity.}
Meiburg's spectral-sensitivity construction uses tournament-indexed rows
of positive conditions and one selected blocking coordinate in each
outgoing row \cite[Sections~3.1 and~5.2--5.5]{Meiburg2026}.\footnote{In the
attribution accompanying this construction, Meiburg credits GPT-5.6-Sol
with the gated-tournament ideas and describes his own contributions as
exposition, size and parameter optimization, and checking the faithfulness
of the Lean formalization.}
Random labels exclude small local configurations in which many accepting
subcubes lie close to the same accepting input. We adapt this architecture
to nested predicates, using a different forbidden-label condition and a
direct union bound. Simultaneous changes of two predicates then control
ordinary sensitivity under recursion.

The distinction between the two sensitivity measures matters here.
The sensitivity graph joins adjacent cube inputs exactly when the function
changes value. Its maximum degree is \(s(f)\); spectral sensitivity
\(\lambda(f)\) is the operator norm of its adjacency matrix
\cite{AaronsonEtAl2021}.
Since \(\lambda(f)\le s(f)\), a lower bound on
\(\bs(f)/\lambda(f)^2\) alone does not supply the ordinary-sensitivity
separation in Theorem~\ref{thm:main}.
The needed estimates are proved directly below, without a spectral input.

\paragraph{Construction and proof strategy.}
We recursively construct nested families of Boolean predicates: their
accepting sets decrease as the predicate index increases. At a recursive
step, the input is partitioned into disjoint child inputs arranged in
rows, and every pair of rows is given one directed edge. Each edge is
labelled by a position in its destination row. A clause associated with
row \(i\) has target conditions requiring the strongest child predicate
to accept every child in that row. For each edge from \(i\) to another row,
it also requires a weaker child predicate to reject the child selected by
the edge label. The parent predicate accepts when one of these clauses
holds. Varying the weaker predicate used in these gate conditions produces
a new nested family.

Fix one parent predicate at a rejecting input, and consider clauses that
can become satisfied after one bit changes. A bit belongs to only one child
and hence affects only one condition of a fixed clause. The candidate
clauses therefore fail exactly one target condition or one gate condition,
and no other condition. The edge labels bound the number of candidates of
the first kind. The tournament allows at most three candidates of the
second kind, and all but at most one have their failed gate in a row whose
strongest child predicates all accept. Repairing such a gate must reverse
both its weaker predicate and the strongest predicate on the same child.
We consequently track the number of bits that reverse a fixed pair of
distinct predicates simultaneously, in addition to ordinary sensitivity.
For a gate repair that makes a clause hold for two rejecting parent
predicates, the two corresponding child gate predicates must both change,
even in the possible exceptional clause.

The initial family consists of indicators of unions of Hamming balls
with successive integer radii around separated centres. One bit cannot
cross two of these radii, so all its joint sensitivities vanish.
Section~\ref{sec:construction} gives this family, the required edge
labelling, and common disjoint sensitive blocks at zero.
Section~\ref{sec:recurrences} then proves the four coupled sensitivity
bounds for arbitrary inputs. With a size parameter \(M\) chosen
sufficiently large for each fixed depth, their normalization in
Section~\ref{sec:separation} bounds sensitivity growth by essentially
a factor \(M\) per level, while the number of disjoint sensitive blocks
grows by \(2M^2+1\). An OR of disjoint copies balances the sensitivities
at zero and one, following the principle used by Ambainis and Sun
\cite[Section~4, Lemma~1]{AmbainisSun2011}.

Finally, Lemma~\ref{lem:composition} amplifies a fixed seed with
\(f(0)=0\) and \(\bs(f,0)>s(f)^2\) into the fixed power separation of
Corollary~\ref{cor:power-separation}. Ambainis and Pr\={u}sis explicitly
noted that a seed with block sensitivity greater than sensitivity squared
would yield a superquadratic power separation by iteration
\cite[Section~1]{AmbainisPrusis2014}. Composition as an amplification method
is treated systematically by Tal~\cite{Tal2013}; we prove the sensitivity
upper bound and zero-input block lower bound used here.
